\section{Balance Sheet Substantiation and Dual Valuation}
\label{sec:substantiation}

The balance sheet is a projection of the move stream, not a separate record. Each move carries CDM-native metadata linking its balance to the originating product, business event, and counterparty. Every line item therefore traces to its business event along a standardised evidence chain. Substantiation is the demonstration to an external auditor that a balance is complete, accurate, and properly valued. Under this construction it reduces to replaying the log. The balance sheet and the ledger are the same stream read two ways, so the two cannot disagree.

\subsection{Balance Reconstruction}

The move stream is an immutable, chronological log, so any balance can be recovered by replaying it. The recipe: start from the opening balance, add every move in, subtract every move out. The balance of wallet $w$ in unit $u$ at time $t$ is its opening balance plus the signed sum of every move on $u$ up to $t$:
\[
w_t(u) = w_0(u) + \sum_{\substack{\text{moves } m:\; m.\text{unit}=u,\; m.\text{time}\leq t}}
\begin{cases}
+m.\text{quantity} & m.\text{destination}=w \\
-m.\text{quantity} & m.\text{source}=w.
\end{cases}
\]
The result is deterministic and reproducible from the raw stream alone. An auditor verifies any balance at any date $t$ by replaying the stream to $t$ and applying this formula. No independent balance record exists to reconcile against.

The reconstruction is a projection of the move stream: each move maps to a two-cell delta, minus at its source and plus at its destination, and the balance map is the sum of these deltas. Pointwise addition, with the all-zero balance as unit, makes the per-$(\text{wallet},\text{unit})$ balance map a \emph{monoid}. The projection \lstinline!balances! is \lstinline!foldMap! over that monoid.

\begin{lstlisting}[language=Haskell]
-- Scalars and keys in local pedagogical form: Qty is the additive abelian
-- group of exact minor units (Section 2); the Move is abridged to four
-- fields with a raw Qty magnitude. The consolidated form -- the reference
-- Move, magnitude PosQty -- is reference/Ledger.hs (Part B).
newtype Qty      = Qty Integer     deriving (Eq, Show)
newtype WalletId = WalletId String deriving (Eq, Ord, Show)
newtype UnitId   = UnitId   String deriving (Eq, Ord, Show)

instance Semigroup Qty where Qty a <> Qty b = Qty (a + b)
instance Monoid    Qty where mempty         = Qty 0
qneg :: Qty -> Qty
qneg (Qty n) = Qty (negate n)

data Move = Move { from, to :: WalletId, unit :: UnitId, qty :: Qty }

-- The projection target: one signed total per (wallet, unit). The monoid
-- is POINTWISE addition -- combine cell by cell, all-zero map as unit. This
-- is NOT the library default for Map, which is left-biased union and would
-- silently drop the second contribution to a cell. Choosing the right
-- monoid is the whole correctness content of this type.
newtype Balances = Balances (Map (WalletId, UnitId) Qty)

instance Semigroup Balances where
  Balances a <> Balances b = Balances (Map.unionWith (<>) a b)
instance Monoid Balances where
  mempty = Balances Map.empty

-- One move debits its source and credits its destination: the two-cell
-- signed contribution the displayed formula writes with + / -.
contribution :: Move -> Balances
contribution m = Balances (Map.fromListWith (<>)
  [ ((to   m, unit m),      qty m)
  , ((from m, unit m), qneg (qty m)) ])

-- foldMap: map every move to its contribution, then combine them all with
-- one associative operation; the empty stream maps to the all-zero balance.
-- Folding the chronological prefix up to t gives w_t; the opening balance
-- is just the moves before the window, so there is no separate opening term
-- and no separate balance record to reconcile against.
balances :: [Move] -> Balances
balances = foldMap contribution

netBal :: WalletId -> UnitId -> [Move] -> Qty
netBal w u ms = let Balances b = balances ms
                in Map.findWithDefault mempty (w, u) b
\end{lstlisting}

\begin{remark}[Balance-as-fold]
Plain reading: balances are not stored and maintained; they are recomputed by folding the move stream. Formally, \lstinline!balances! is a \emph{monoid homomorphism} from the free monoid of move lists (lists under concatenation, empty list as unit) into the monoid of balance maps---it carries concatenation to combination,
\[
\texttt{balances (xs ++ ys) = balances xs <> balances ys}.
\]
This one law is the substantiation guarantee read three ways. With \texttt{ys} empty and the operands reassociated, the grouping of independent batches cannot change the result, so replay is order-stable and deterministic. With \texttt{xs} the moves up to a snapshot and \texttt{ys} the moves since, the snapshot balance combined with later moves equals the balance of the whole stream---the $O(k)$ snapshot query of Section~\ref{subsec:snapshots} is correct \emph{by this law}, not by separate argument. And summing any unit's column over all wallets cancels each $+q$ against its $-q$ to \texttt{mempty}: that is the conservation law P1 (Section~\ref{sec:ledger}), here a consequence of \lstinline!contribution! emitting equal and opposite legs, not an assertion to enforce.
\end{remark}

\subsection{Substantiation Scope}

Position quantities and transaction provenance are substantiated by construction. The position set, the flows between wallets, and unit state at any historical date all derive from the move stream by the reconstruction above. This eliminates the manual comparison of independent records that substantiation otherwise requires.

Valuation and disclosure are not so substantiated. Fair value hierarchy classification under IFRS~13 (Level~1/2/3, with sensitivity analysis for Level~3 inputs) depends on inputs outside the move stream. So do the disclosures of IAS~1 paragraphs 117--124 (accounting policies, significant estimates and judgments). Those inputs are external price feeds, valuation-model documentation, management judgment, and audit assessment. Legal enforceability depends on external documentation. The design claim is therefore exact: the ledger substantiates quantities and transaction provenance by construction; valuations, disclosures, and legal status require supplementary evidence.

Unit state read during substantiation comes from UnitStatus, a projection of the same log (\S\ref{sec:states-unitstatus}).

Appendix~\ref{app:reconciliation} enumerates, in one place, the internal reconciliation failure modes the construction renders structurally unreachable, the boundary modes that remain detectable, and the by-design treatment of securities lending.

\subsection{Dual Valuation}

No single price serves settlement and risk management at once. Settlement requires the price at which an external party will transact. Risk management requires a single coherent price across related positions. The resolution is to carry two prices per unit while keeping one position set: both views read the same positions, and only the price vector differs.

\begin{mentalmodel}
Picture each position as water sealed in a fixed run of pipes, with two rate sheets posted on the outside of the seal, each converting the one water level into money by its own rule. Settlement values the water off one sheet and risk off the other; because both read the same water, the quantity beneath them can never diverge. The two money figures do diverge, and either can move while not a drop of water shifts --- unlike the sealed water, neither figure is conserved, for the sheets hang outside the seal.
\end{mentalmodel}

\begin{definition}[Dual Valuation]
Each unit $u\in\mathcal{U}$ carries two price functions:
\begin{itemize}[nosep]
\item $P_t^{\text{MtMk}}(u)$, mark-to-market, from observable market data---used for settlement, margin, external reporting, and custodian reconciliation;
\item $P_t^{\text{MtMd}}(u)$, mark-to-model, from internal valuation models---used for risk management, portfolio optimisation, and performance attribution.
\end{itemize}
Their differential is $\Delta_t(u)=P_t^{\text{MtMd}}(u)-P_t^{\text{MtMk}}(u)$.
\end{definition}

\paragraph{Multi-exchange futures.} Nikkei~225 futures settle on three exchanges at different cut-off times:
\begin{align*}
P_t^{\text{MtMd}}(\text{NKY}) &= 35{,}500 && \text{(theoretical fair value)}\\
P_t^{\text{MtMk}}(\text{NKY-SIMEX}) &= 35{,}480 && \text{(Singapore settlement)}\\
P_t^{\text{MtMk}}(\text{NKY-CME}) &= 35{,}490 && \text{(Chicago settlement)}\\
P_t^{\text{MtMk}}(\text{NKY-OSE}) &= 35{,}495 && \text{(Osaka settlement)}.
\end{align*}
The three cannot share a mark-to-market price. Singapore closes hours before Chicago, and Osaka closes between them. Each settlement price fixes the market at a different moment. Margin must therefore use the exchange-specific price. They share one mark-to-model price because all converge to the same final settlement value, the Nikkei~225 index. A desk holding offsetting positions across the three reads net exposure from the single MtMd price. Each leg settles at its own MtMk price.

\paragraph{Fair value adjustment.} The model price admits an adjustment for the gap between theoretical and executable prices: liquidity cost, bid-ask spread, model uncertainty. The adjusted price is
\[
P_t^{\text{MtMd, adj}}(u) = P_t^{\text{MtMd}}(u) + \text{FVA}_t(u).
\]
Prudent valuation under CRR Article~105 requires the adjusted price for regulatory capital. Risk management uses the unadjusted price. Dual valuation carries both.

\paragraph{Single position set.} Desk-level consistency and firm-level reconciliation coexist because both valuations read the same wallet balances. Each desk prices related positions coherently through $P^{\text{MtMd}}$. The firm reads valuation discrepancy directly from $\Delta_t(u)$. Internal transfers may book on either basis by policy. Position divergence between the two views is impossible: there is one position set and two price vectors.

\subsection{Snapshots and Audit Trail}
\label{subsec:snapshots}

A snapshot is a substantiated checkpoint: a consistent cut of the move stream, carrying both pricing bases and the provenance metadata of each move. A balance query starting from the most recent snapshot replays only the moves since it. The query is therefore $O(k)$ in the moves since the snapshot rather than $O(n)$ in the total. A snapshot is consistent only when taken at a quiescence point (no transactions in flight) or under multi-version concurrency control.

Quantities and prices separate cleanly: the quantity layer is closed and constrained by conservation, and prices enter only afterwards as a multiplier. Conservation enforces double-entry accounting on quantities, independent of price:
\[
\sum_{w\in\mathcal{W}} w_t(u) = 0 \qquad \forall u,\,\forall t,
\]
the closed-system law of Section~\ref{sec:ledger} (P1). Valuation applies the price vector to the conserved quantities:
\[
V_t^{\text{MtMk}} = \sum_{w} w_t(u)\cdot P_t^{\text{MtMk}}(u), \qquad
V_t^{\text{MtMd}} = \sum_{w} w_t(u)\cdot P_t^{\text{MtMd}}(u).
\]
Quantities are conserved and shared. Only the two valuations differ.

